% ECONOMICS_PROBLEM: {"id": "C90-642", "title": "External validity of experimental effects", "jel_code": "C90", "source_id": "OP-0061"}
% !TeX program = xelatex
\documentclass[11pt,a4paper]{article}
\usepackage[margin=23mm]{geometry}
\usepackage{fontspec}
\setmainfont{Latin Modern Roman}
\usepackage{amsmath,amssymb,mathtools}
\usepackage{xcolor,enumitem,booktabs,longtable,array,xurl,hyperref}
\definecolor{muted}{HTML}{53636C}
\hypersetup{colorlinks=true,urlcolor=blue,linkcolor=blue}
\setlength{\parindent}{0pt}
\setlength{\parskip}{5pt}
\newcommand{\R}{\mathbb R}
\newcommand{\E}{\mathbb E}
\newcommand{\N}{\mathbb N}
\newcommand{\ind}{\mathbf 1}
\newcommand{\Var}{\operatorname{Var}}
\newcommand{\Cov}{\operatorname{Cov}}
\newcommand{\supp}{\operatorname{supp}}
\DeclareMathOperator*{\argmin}{arg\,min}
\DeclareMathOperator*{\argmax}{arg\,max}
\newcommand{\entrymeta}[3]{{\sffamily\small\color{muted}\textbf{\texttt{#1}}\quad|\quad #2\par #3\par}}
\newcommand{\Problemref}[1]{\texttt{#1}}
\newcommand{\JELref}[2]{\texttt{#2} (JEL \texttt{#1})}
\begin{document}
\section*{External validity of experimental effects}
\textbf{Problem ID:} \texttt{C90-642}\quad \textbf{JEL:} \texttt{C90}\par
% BEGIN ECONOMICS STATEMENT
\label{problem:OP-0061}
\label{atlas:prob:B1}

\noindent\textbf{Original identifier:} B1 (atlas item 61). \textbf{Classification:} Causal transportability research program. \textbf{Original tractability label:} Medium (editorial, not a complexity result).

\subsection*{Definitions and assumptions}

Let $S=1$ denote participation in a source experiment and $S=0$ membership in a specified target population. Let $X$ be pretreatment covariates, $A\in\{0,1\}$ randomized treatment, and $c\in\mathcal C$ an institutional context. The potential outcome $Y(a,c)$ is measured on a common, prespecified scale. Assume consistency and no interference within each context; interference requires a different estimand. Randomization identifies source conditional effects wherever treatment has positive probability. Sampling positivity requires every target covariate profile to occur with positive source probability. Distinguish selection into the sample from changing the intervention itself.

\subsection*{Formal research question}

For fixed source and target contexts $c_s,c_t$, identify or sharply bound
\[
\tau_t=\mathbb E[Y(1,c_t)-Y(0,c_t)\mid S=0].
\]
Determine substantively testable conditions under which
\[
\tau_t=\int \mathbb E[Y(1,c_s)-Y(0,c_s)\mid X=x,S=1]\,dF_{X\mid S=0}(x).
\]
This equality follows from conditional equality of treatment effects across populations and contexts, together with overlap; randomization alone does not imply it. When equality is relaxed by a specified discrepancy bound $b(x)$, derive target-effect bounds and confidence regions accounting for source estimation and target sampling.

\subsection*{Interpretation and scope}

The central extension is a prospectively validated transport rule for a defined mechanism, outcome, and target population. Paired experiments should hold the treatment protocol constant where possible and independently vary stakes, scrutiny, experience, and institutional rules. Reweighting removes measured compositional differences only under the stated conditional effect restriction. It cannot automatically remove changes in incentives or unmeasured participation motives. A multisite model must distinguish site-level heterogeneity from individual-level sampling error, and predictive coverage should be evaluated on held-out contexts. A finding that the average effect transports need not establish transport of distributional effects or welfare. The relevant failure criterion is an error exceeding a prespecified policy tolerance, rather than a failure to reject equal effects in a small sample.

\subsection*{Source audit and status}

All three author-year references are verified. [S1] supplies a field-context taxonomy; [S2] discusses generalization threats. [S3] concerns replication of laboratory studies and is only indirectly relevant: replicability does not establish external validity. None states the displayed transport equation as a published open conjecture. This is an editorial identification and validation problem motivated by these sources; its unresolved scope is conditional on the selected target and model restrictions, not certified for every application in 2026.

\subsection*{References}

\begin{enumerate}[label={[S\arabic*]},leftmargin=*]

\item Glenn W. Harrison and John A. List (2004). \emph{Field Experiments}. Journal of Economic Literature 42(4), 1009--1055. \url{https://www.aeaweb.org/articles?id=10.1257/0022051043004577}\par \textit{Source role:} Taxonomy of field settings; methodological foundation. \textit{Locator:} Abstract; six dimensions of field context.

\item Steven D. Levitt and John A. List (2007). \emph{What Do Laboratory Experiments Measuring Social Preferences Reveal About the Real World?}. Journal of Economic Perspectives 21(2), 153--174. \url{https://www.aeaweb.org/articles?id=10.1257/jep.21.2.153}\par \textit{Source role:} Methodological discussion of generalizability. \textit{Locator:} Abstract; discussion of scrutiny, context, selection, and stakes.

\item Colin F. Camerer, Anna Dreber, Eskil Forsell, Teck-Hua Ho, Jürgen Huber, Magnus Johannesson, Michael Kirchler, Johan Almenberg, Adam Altmejd, Taizan Chan, Emma Heikensten, Felix Holzmeister, Taisuke Imai, Siri Isaksson, Gideon Nave, Thomas Pfeiffer, Michael Razen, and Hang Wu (2016). \emph{Evaluating Replicability of Laboratory Experiments in Economics}. Science 351(6280), 1433--1436. \url{https://doi.org/10.1126/science.aaf0918}\par \textit{Source role:} Evidence on laboratory replicability, not field transportability. \textit{Locator:} Abstract; author repository metadata and supplemental materials.

\end{enumerate}

\subsection*{Common conventions: Source mathematical conventions}

Unless an entry states otherwise, agent and firm sets are finite. State and action spaces are measurable, strategies and outcome maps are measurable, probability laws are specified, and expectations exist for the targets under discussion. Infinite-horizon discount factors lie in $(0,1)$ and discounted objectives are finite. Norms, tolerances and complexity input sizes must be specified for computational comparisons. Constants or primitives that appear locally are inputs, not empirical facts asserted by this catalogue.

\textbf{Identification.} A statistical model $m\in\mathcal M$ implies an observable law $P_m$ and target $\theta(m)$. At law $P$, its identified set is
\[
\Theta(P;\mathcal M)=\{\theta(m):m\in\mathcal M,\ P_m=P\}.
\]
Point identification holds when this set is a singleton, or equivalently when $P_{m_1}=P_{m_2}$ implies $\theta(m_1)=\theta(m_2)$ for all compatible models. Sharp bounds mean bounds attained, or approached when the boundary is not attained, by compatible models. Writing this set is a definition, not a solution: the substantive task is to establish informative restrictions and derive or estimate the set. An unrestricted-confounding model may give only trivial bounds.

\textbf{Inference.} Identification, estimability and valid uncertainty quantification are separate questions. Fix $\alpha\in(0,1)$. A confidence procedure $C_n$ over a declared law class $\mathcal P$ has asymptotic uniform coverage at level $1-\alpha$ when
\[
\liminf_{n\to\infty}\inf_{P\in\mathcal P}\Pr_P\{\theta(P)\in C_n\}\geq1-\alpha.
\]
For set-identified targets the coverage event must be defined separately, for example coverage of the full identified set. If an entry uses identification-indexed models instead of uniquely indexed laws, the same coverage convention is applied over those models. Pointwise limits do not establish uniform validity, and asymptotic validity does not guarantee finite-sample precision. Sample sizes, dependence structures and weak-identification sequences are part of each application.

\textbf{Counterfactuals.} $Y_i(a)$ is a potential outcome under a specified intervention, including its timing, implementation and versions. It is not $Y_i$ conditional on voluntarily choosing $a$. A full intervention vector permits interference; shorthand scalar treatment notation does not silently impose no interference. Assignment, selection, attrition, anticipation and measurement must be specified. A claim about an instrument's local effect is not automatically a population policy effect.

\textbf{Economic equilibrium.} A counterfactual model includes best responses, constraints, feasibility, market clearing, state transitions and expectations consistent with the stated information. When several equilibria exist, either a declared selector is an input or the target is set-valued. Comparisons across policy regimes must use the stated selection convention. Reduced-form relationships are not presumed invariant after an equilibrium-changing intervention.

\textbf{Optimization and welfare.} Optimization is over the stated feasible set. If attainment is not established, $\sup$ or $\inf$ and $\varepsilon$-optimal choices replace an asserted maximizer or minimizer. For example, with value $V=\sup_{a\in\mathcal A}W(a)<\infty$, an $\varepsilon$-optimal set is $\{a\in\mathcal A:W(a)\geq V-\varepsilon\}$ for $\varepsilon>0$. Benefits, costs, risk and health or environmental outcomes can be aggregated only using declared units and conversion weights. Interpersonal utility weights, the treatment of fiscal transfers and foreign profits, discounting, and the behavioral welfare criterion are explicit inputs. A comparative-static derivative additionally requires existence and appropriate differentiability of the selected counterfactual.

\textbf{Sources.} Local labels [S1], [S2], and so forth restart in every question. Locators identify the relevant abstract, model, section or result at the level actually checked; a bibliographic or abstract-level verification is not represented as inspection of an entire proof. Working papers are distinguished from later publications and changing versions. Source summaries are paraphrased. All URL checks and editorial formulations refer to the audit date stated above.
% END ECONOMICS STATEMENT
\end{document}
